\documentclass[12pt,a4paper]{article}
\usepackage[utf8]{inputenc}
\usepackage{amsmath, graphicx, hyperref, listings, booktabs}
\usepackage[margin=2.5cm]{geometry}
\title{Assignment 3: The Swiss Cheese Test Suite}
\author{Hikmatullah Danish \quad 3125999089}
\date{July 2026}

\begin{document}
\maketitle

\section{Introduction}
This assignment applies the Swiss Cheese model of accident causation to software testing. The principle: no single test layer catches all errors, but multiple overlapping layers collectively prevent defects from reaching production. A 4-layer test suite was designed for a flood inundation model, with each layer catching what the previous layer missed. The assignment also required identifying at least one AI hallucination—code that passes syntax checks but violates physical reality.

\section{The 4-Layer Model}
\textbf{Layer 1 — Unit Tests (6 tests):} Verifies basic execution: does the function return a dict? Are the required keys present? Is the mask boolean? Are shapes consistent? This layer catches crashes and type errors.

\textbf{Layer 2 — Boundary Tests (5 tests):} Tests extreme values: water level below minimum elevation (0\% flood), above maximum elevation (100\% flood), exactly at boundaries, and with very small increments. This layer catches off-by-one errors and edge-case mishandling.

\textbf{Layer 3 — Physical Tests (10 tests):} Validates scientific constraints: monotonic increase with water level, non-negative depths everywhere, max depth equals water-level minus minimum elevation, non-flooded areas have zero depth, and mask-depth array consistency. This layer catches physically impossible outputs.

\textbf{Layer 4 — Regression Tests (3 tests):} Ensures reproducibility: same seed produces identical DEM, different seed produces different DEM, and deterministic flood calculation. This layer catches non-deterministic behaviour.

\section{Hallucination Detection}
The test suite was designed to catch four known AI weaknesses: (1) computing depth as $W - E_{ij}$ without checking whether it is negative, (2) using integer division for percentage calculation, (3) using $\leq$ instead of $<$ for the flood condition (changing behaviour at exact boundary), and (4) off-by-one errors in array indexing. The boundary layer's test at exact maximum elevation exposed the $<$ vs $\leq$ distinction: with strict inequality, cells exactly at water level are not flooded—a design choice that must be documented.

\section{Results}
All 24 tests pass across the four layers (6 + 5 + 10 + 3). No active hallucinations were detected in the current flood-inundation implementation—the code correctly uses \texttt{np.where()} for depth calculation, strict inequality for flood condition, and float64 for percentage computation. The test suite is self-contained and requires no imports from external experiment folders.

\section{Conclusion}
The Swiss Cheese model provided a systematic framework for covering blind spots that a flat test list would miss. Layer 3 (physical tests) proved most valuable, catching the qualitative difference between syntactic correctness and physical validity. This assignment demonstrates that multi-layer testing is essential for AI-generated scientific code, where each layer catches qualitatively different error classes.

\end{document}